\subsection{A necessary score-order bridge from regularisation to AUC}
\label{sec:v15-score-bridge}
An alternative explanation for a large numerical shrinkage difference
is that two discriminants merely rescale the same test scores.
A positive affine score transformation does not change ROC--AUC.
We therefore audited the untouched Shin2017A \emph{frozen
trial-score arrays} rather than relying on the fitted $\lambda$
values alone.

\paragraph{Proposition 3 (a label-free AUC-change upper bound).}
For $n_0$ negative and $n_1$ positive examples, let
$K_{\mathrm{all}}(s_D,s_C)$ be the number of \emph{all unordered trial
pairs} whose pairwise order differs between score arrays $s_D$ and
$s_C$, with ties treated as a separate order. For any fixed unknown
assignment of binary labels with these class counts,
\[
 \left|\mathrm{AUC}(s_D)-\mathrm{AUC}(s_C)\right|
 \le\min\left\{1,\frac{K_{\mathrm{all}}(s_D,s_C)}{n_0n_1}\right\}.
\]
In particular, if $K_{\mathrm{all}}=0$, the two methods have exactly
the same ROC--AUC under \emph{every} labeling.

\emph{Proof.}
ROC--AUC averages the pairwise ordering indicator over the
$n_0n_1$ between-class pairs, assigning $1/2$ to ties.
Each between-class pair contributes the same amount under
both score arrays unless its pairwise order changes; a
changed pair alters that contribution by at most one.
The changed between-class pairs are a subset of the
$K_{\mathrm{all}}$ changed pairs computable without labels.
The triangle inequality and $0\le\mathrm{AUC}\le1$
give the stated bound.\hfill$\square$

Using the original checksum-locked Shin2017A score arrays
($N=29$, 60 trials per participant, 30 per class), the
source-macro-selected versus covariance-NLL-selected
discriminants changed the order of, on average,
\textbf{13.87\%} of all within-participant trial pairs
(range 3.95--30.23\%). The median trial-score Spearman
correlation was 0.894. Every one of the 29 published
held-person AUC contrasts satisfies the label-free
necessary upper bound, and no participant has
identical trial-score rankings across the two selections.
Thus the changed fitted shrinkage region is not
\emph{merely} a positive score rescaling on the observed
trials.

The inversion count is a \emph{necessary consistency
mechanism}, not a sufficient predictor of gain, since
only between-class ordering enters AUC and the
label-free count also includes within-class pairs.
It is not a newly tuned classifier, a new prospective
experiment or evidence that target-specific class
structure can be identified without labels. The
29 per-person comparisons share substantial source
training data and must not be treated as independent
mechanistic replications. The verification script
independently reads the original ZIP and fails if
any reported AUC change exceeds its mathematically
permitted ordering bound.
